\documentclass[10pt,a4paper]{article}
\usepackage{amsmath,amssymb,amsthm,graphicx}
\usepackage{amsmath,amssymb,amsthm,mathtools,bm,enumerate}
\usepackage{hyperref}
\usepackage{cleveref}
\newtheorem{theorem}{Theorem}
\newtheorem{definition}[theorem]{Definition}
\newtheorem{lemma}[theorem]{Lemma}
\newtheorem{claim}[theorem]{Claim}
\newtheorem{corollary}[theorem]{Corollary}
\newtheorem{remark}[theorem]{Remark}
\newtheorem{THM}[theorem]{Theorem}
\crefname{THM}{theorem}{theorems}
\Crefname{THM}{Theorem}{Theorems}
\crefname{claim}{claim}{claims}
\Crefname{claim}{Claim}{Claims}
\DeclareMathOperator*{\Ex}{\mathbb E}
\DeclareMathOperator{\Supp}{Supp}
\newcommand{\ex}[2][]{\Ex_{#1}\left[#2\right]}
\newcommand{\pr}[2][]{\Pr_{#1}\left[#2\right]}
\newcommand{\abs}[1]{\left|#1\right|}
\newcommand{\inp}[1]{\left\langle#1\right\rangle}
\newcommand{\norm}[2][]{\left\|#2\right\|_{#1}}
\DeclarePairedDelimiter\ceil{\lceil}{\rceil}
\DeclarePairedDelimiter\floor{\lfloor}{\rfloor}
\newcommand{\bbR}{\mathbb R}
\newcommand{\bbN}{\mathbb N}
\newcommand{\cA}{\mathcal A}
\newcommand{\cB}{\mathcal B}
\newcommand{\cC}{\mathcal C}
\newcommand{\cE}{\mathcal E}
\newcommand{\cG}{\mathcal G}
\newcommand{\bits}[1]{\{0,1\}^{#1}}
\newcommand{\bit}{\{0,1\}}
\newcommand{\eps}{\varepsilon}
\newcommand{\wt}{\mu}
\newcommand{\poly}{\mathrm{poly}}
\newcommand{\sout}[1][]{s_{\mathrm{out}#1}}
\newcommand{\si}[1][]{s_{\mathrm{in}#1}}
\newcommand{\Si}{S_{\mathrm{in}}}
\newcommand{\ninf}{\norm}
\newcommand{\nmax}{\norm[\mathrm{max}]}
\newcommand{\mathbbm}[1]{\mathbf{#1}}
\allowdisplaybreaks
\newcommand{\BPL}{\mathbf{BPL}}
\renewcommand{\L}{\mathbf{L}}
\newcommand{\zo}{\{0,1\}}
\title{Optimal-error signed pseudodistributions for read-once branching programs}
\author{Eshan Chattopadhyay, Jyun-Jie Liao}
\date{doi:10.4230/LIPIcs.CCC.2020.25; CC BY 3.0 Unported}
\begin{document}
\maketitle
\noindent\textit{Editorial note.} Editable transcription only of formally published CCBY3.0Unported mainPRPD result andfullbearingproofs. Nonexclusive prepublicationnativev4 private aid; publicselected filesexclude allunselectedranges/comments. EntireSections2–6 andAppendixB retained; AppendixA SaksZhou sketch excluded. Originalauthorparameters, norms, signedweights, indices anderrorrecurrences retained without mathematicalrepair oraddedhypotheses. PRPD expectation uses signedrho andis not an ordinary probability distribution; noBPL=L claim. Originalglobal theorem numbering restored andcustomTHM reference type registered. Ordinarybold1 indicator replaces unavailablebbmfont only. Noinput/include/assets. Allproduction/validationYupeng; rootownscompilation.\par\medskip
\noindent\textit{Allowed standard prerequisites (editorial).}\begin{itemize}
\item ROBP transition matrices, averaging samplers/extractors and infinity/max matrix norms. Complete formalSections2–6 allsource-local matrix/robustapproximation/recursive proofs; citedstandardextractor facts andcompleteAppendixB sampler construction retained
\end{itemize}\medskip
\begin{definition}[Read-once branching program]
A $(n,w)$-read-once branching program (ROBP) $B$ is a directed graph on the vertex set $V=\bigcup_{i=0}^n V_i$, where each set $V_i$ contains $w$ nodes. Every edge in this directed graph is labeled either $0$ or $1$. For every $i<n$, and every node $v\in V_i$, there exists exactly two edges starting from $v$, one with label $0$ and the other with label $1$. Every edge starting from a node in $V_i$ connects to a node in $V_{i+1}$.
We say $n$ is the \emph{length} of $B$, $w$ is the \emph{width} of $B$ and $V_i$ is the $i$-th \emph{layer} of $B$.

Moreover, there exists exactly one starting state $s\in V_0$, and exactly one accepting state $t\in V_n$. For every $x=(x_1,\dots,x_n)\in\bits{n}$, we define $B(x)=1$ if starting from $s$ we will reach $t$ following the edges labeled by $x_1,\dots,x_n$. Otherwise we define $B(x)=0$.
\end{definition}
It is well-known the computation of a probabilistic Turing machine that uses space $S$ and tosses $n$ coins,  on a given input $y$, can be carried out by a $(n,2^{O(S)})$-ROBP $B_y$. In particular, if the string $x \in \zo^n$ corresponds to the $n$ coin tosses, then  $B_y(x)$ is the output of the Turing machine. 

A standard derandomization technique is via \emph{pseudorandom generators}. We define this notion for the class of ROBPs.

\begin{definition}[Pseudorandom generator]
A function $G:\bits{s}\to\bits{n}$ is a $(n,w,\eps)$-pseudorandom generator (PRG) if for every $(n,w)$-ROBP $B$,
$$\abs{\ex[x\in\bits{n}]{B(x)}-\ex[r\in\bits{s}]{B(G(r))}}\le\eps.$$
The \emph{seed length} of $G$ is $s$.   $G$ is explicit if $G$ is computable in $O(s)$ space.
\end{definition}
To derandomimze  space-bounded computation  given an explicit $(n,w,\eps)$-PRG, one can enumerate $B(G(r))$ for every $r\in\bits{s}$ with $O(s)$ additional space to compute an $\eps$-approximation of the quantity $\ex[x]{B(x)}$. 

Nisan~\cite{Nis92} constructed a $(n,w,\eps)$-PRG with seed length $O(\log n \cdot \log (nw/\eps))$, which implies $\BPL\subseteq \L^2$.   While there is a lot of progress in constructing PRG with better seed length for restricted family of ROBP (see, e.g., \cite{NZ96,Arm98,BV10,BDVY13,BRRY14,KNP11,De11,Ste12,MRT19} and references therein), Nisan's generator and its variants~\cite{Nis92,INW94,RR99} remain the best-known generators in the general case.
\subsection{Pseudorandom pseudodistribution }
Recently, a beautiful work of Braverman, Cohen and Garg \cite{BCG18} introduced the notion of a \emph{pseudorandom pseudodistribution} (PRPD)  that relaxes the definition of a PRG.
\begin{definition}[Pseudorandom pseudodistribution]
A pair of functions $(G,\rho):\bits{s}\to \bits{n}\times \bbR$ generates a $(n,w,\eps)$-pseudorandom pseudodistribution (PRPD) if for every $(n,w)$-ROBP $B$,
$$\abs{\ex[x\in\bits{n}]{B(x)}-\ex[r\in\bits{s}]{\rho(r)\cdot B(G(r))}}\le\eps.$$
We say $s$ is the \emph{seed length} of $(G,\rho)$. We say $(G,\rho)$ is $k$-bounded if $|\rho(x)|\le k$ for every $x\in\bits{s}$. We say $(G,\rho)$ is explicit if they are computable in space $O(s)$.
\end{definition}

\begin{THM}
\label{thm:main}
There exists an explicit $(n,w,\eps)$-PRPD generator $(G,\rho)$ with seed length
$$O\left(\log n\cdot \log(nw)\cdot \log\log (nw) + \log(1/\eps)\right),$$
which is $\poly(1/\eps)$-bounded.
\end{THM}

\section{ROBPs and Matrices}\label{sec:matrices}
We introduce the matrix representation of ROBPs and some related definitions that are useful in the rest of the paper. First, we setup some notation.
\newline
\textbf{Notation:} Given two strings $x,y$, we use $x\Vert y$ to denote the concatenation of $x$ and $y$. For every $n\in\bbN$, we use $[n]$ to denote the set $\{1,2,\dots,n\}$. We denote a collection of objects $A_i^j$ with subscript $i\in S$ and superscript $j\in T$ by $[A]_S^T$ for short.

Given a $(n,w)$-ROBP $B$ with layers $V_0,\dots,V_n$, we can represent the transition from layer $V_{t-1}$ to $V_{t}$ by two stochastic matrices $M_t^0$ and $M_t^1$ as follows: suppose layer $V_j$ consists of the nodes $\{v_{j,1},\dots,v_{j,w}\}$.  The entry $(M_t^0)_{i,j}=1$ if and only if there exist a $0$-labeled edge from $v_{t-1,i}$ to $v_{t,j}$ (else $(M_t^0)_{i,j}=0$). The matrix $M_t^1$ is defined similarly according to the edges that labeled $1$ between layers $V_{t-1}$ and $V_t$. More generally, we can also represents multi-step transition by a stochastic matrix. That is, for every $0\le a\le b\le n$, and every $r=(r_{a+1},\ldots,r_b)\in\bits{b-a}$, we can define
$$M_{a..b}^r=\prod_{t=a+1}^b M_t^{r_t}$$
which corresponds to the transition matrix from layer $a$ to layer $b$ following the path labeled by $r$. Note that every row of $M_{a,b}^r$ contains exactly one $1$, and the other entries are $0$.

An $n$-step random walk starting from the first layer can be represented with the following matrix:
$$M_{0..n}=\frac{1}{2^n}\sum_{r\in\bits{n}}M_{0..n}^r=\prod_{t=1}^n \frac{1}{2}\left(M_t^0+M_t^1\right).$$
By definition of $M_t^0,M_t^1$ one can observe that the $(i,j)$ entry of $M_{0..n}$ is the probability that a random walk from $v_{0,i}\in V_0$  reaches $v_{n,j}\in V_n$. Therefore, suppose $v_{0,i}\in V_0$ is the starting state of $B$, $v_{n,j}\in V_n$ is the accepting state of $B$, then $\ex[x]{B(x)}$ equals the $(i,j)$ entry of $M_{0..n}$. 

Recall that a generator of a $(n,w,\eps)$-PRPD is a pair of function $(G,\rho)$ such that for every $(n,w)$-ROBP $B$,
$$\abs{\ex[r]{\rho(r)\cdot B(G(r))}-\ex[x\in\bits{n}]{B(x)}}\le \eps.$$
Equivalently, for every transition matrices $M_1^0,M_1^1,\ldots,M_n^0,M_n^1$, we have 
$$\nmax{\ex[r]{\rho(r)\cdot M_{0..n}^{G(r)}} - M_{0..n}}\le \eps,$$
where $\nmax{A}$ denotes $\max_{i,j}\abs{A(i,j)}$. 

Therefore it is natural to represents a PRPD $(G,\rho)$ with a mapping $\cG:\bits{s}\to \bbR^{w\times w}$ where $\cG(r)=\rho(r)\cdot M^{G(r)}_{0..n}$. More generally, we will use a notation similar to the ``matrix bundle sequence" (MBS) introduced in \cite{BCG18} to represent a PRPD.

\begin{definition}\label{def:MBS}
Consider a $(n,w)$-ROBP $[M]_{[n]}^{\bit}$ and a pair of functions $(G,\rho):\bits{\sout}\times[\Si]\to\bits{n}\times\bbR$. The \emph{matrix form} of $(G,\rho)$ on $[M]_{[n]}^{\bit}$ is a mapping $\cA:\bits{\sout}\times[\Si]\to\bbR^{w\times w}$ such that for every $x\in\bits{\sout}$ and $y\in[\Si]$,
$$\cA(x,y)=\rho(x,y)\cdot M_{0..n}^{G(x,y)}.$$
For every $x\in\bits{\sout}$ we abuse the notation and define
$$\cA(x)=\ex[y]{\cA(x,y)}.$$
Besides, we define $\inp{\cA}=\ex[x,y]{\cA(x,y)}$. We say $\sout$ is the \emph{outer seed length} of $\cA$, denoted by $\sout(\cA)$, and $\Si$ is the \emph{inner size} of $\cA$, denoted by $\Si(\cA)$. We also define $\si(\cA)=\ceil{\log \Si}$ to be the \emph{inner seed length} of $\cA$, and $s(\cA)=\sout(\cA)+\si(\cA)$ to be the \emph{seed length} of $\cA$.
\end{definition}
\begin{remark}
For every fixed $x$, the collection $\{\cA(x,y):y\in[\Si]\}$ corresponds to the ``matrix bundle" in \cite{BCG18}. This should be treated as a collection of matrices which ``realizes" the matrix $\cA(x)$. The whole structure $\cA$ corresponds to the ``matrix bundle sequence" in \cite{BCG18}, and should be treated as a uniform distribution over the set $\{\cA(x):x\in\bits{\sout}\}$. 
\end{remark}

When the ROBP $[M]_{[n]}^{\bit}$ is clear in the context, we will use the matrix form $\cA$ to represent the pseudodistribution $(G,\rho)$ directly. We will apply arithmetic operations on matrices $\cA(x)$, and these operations can be easily translated back to operations on pseudodistributions as follows.
\begin{definition}
Consider a $(n,w)$-ROBP $[M]_{[n]}^{\bit}$, and a pair of function $(F,\sigma):[S]\to \bits{n}\times\bbR$. The matrix that is \emph{realized} by $(F,\sigma)$ on $M_{0..n}$ is $\ex[{i\in[S]}]{\sigma(i)\cdot M_{0..n}^{F(i)}}$. We say S is the \emph{size} of $(F,\sigma)$.
\end{definition}
\noindent
Scaling the matrix corresponds to scaling the coefficients in the pseudodistribution.
\begin{claim}
Consider a $(n,w)$-ROBP $[M]_{[n]}^{\bit}$, let $A$ be a matrix realized by matrix bundle $(F_A,\sigma_A)$ on $M_{0..n}$. Then $cA$ is realized by a matrix bundle $(F'_A,\sigma'_A)$ of size $S_A$ s.t. $F'_A=F_A$ and $\sigma'_A(x)=c\sigma_A(x)$ for every $x\in[S]$.
\end{claim}
\noindent
The summation on matrices corresponds to re-weighting and union on pseudodistributions.
\begin{claim}
Consider a $(n,w)$-ROBP $[M]_{[n]}^{\bit}$, let $A$ be a matrix realized by matrix bundle $(F_A,\sigma_A)$ of size $S_A$ on $M_{0..n}$ and $B$ be a matrix realized by matrix bundle $(F_B,\sigma_B)$ of size $S_B$ on $M_{0..n}$. Then $A+B$ is realized by a matrix bundle $(F',\sigma')$ of size $S_A+S_B$ on $M_{0..n}$ s.t. 
\[
F'(x) = \begin{cases} F_A(x) &\mbox{if } x \le S_A \\
F_B(x-S_A) & \mbox{if }  x > S_A \end{cases}\textrm{ and }
\sigma'(x) = \begin{cases} \frac{S_A+S_B}{S_B}\cdot \sigma_A(x) &\mbox{if } x \le S_A \\
\frac{S_A+S_B}{S_B}\cdot \sigma_B(x-S_A) & \mbox{if }  x > S_A \end{cases}
\]
\end{claim}
\noindent
The multiplication on matrices corresponds to concatenation of pseudodistributions.
\begin{claim}
Consider a $(n,w)$-ROBP $[M]_{[n]}^{\bit}$, let $A$ be a matrix realized by matrix bundle $(F_A,\sigma_A)$ of size $S_A$ on $M_{0..n/2}$ and $B$ be a matrix realized by matrix bundle $(F_B,\sigma_B)$ of size $S_B$ on $M_{n/2..n}$. Fix a bijection $\pi:[S_A]\times[S_B]\to{[S_A\cdot S_B]}$. Then $AB$ is realized by a matrix bundle $(F',\sigma')$ of size $S_A\cdot S_B$ s.t. for every $a\in[S_A],b\in[S_B]$,
$$F'(\pi(a,b))=F_A(a)\Vert F_B(b)\textrm{ and }\sigma'(\pi(a,b))=\sigma(a)\cdot \sigma(b).$$
\end{claim}

\section{Proof Overview}\label{sec:proof-overview}
In this section we give an outline of our construction and proof. In \Cref{sec:samp-arg}, we briefly recap how a sampler is used in \cite{BCG18} to achieve better seed length in the small-error regime. We discuss our construction ideas in \Cref{sec:our-const}. 
\subsection{The sampler argument}\label{sec:samp-arg}
 Nisan's generator and its variants recursively use a lemma of the following form.
\begin{lemma}\label{lemma:INW}
Consider a $(n,w)$-ROBP $[M]_{[n]}^{\bit}$. Let $\cA$ be the matrix form of a distribution on $M_{0..n/2}$, and $\cB$ be the matrix form of a distribution on $M_{n/2..n}$. Suppose $s(\cA)=s(\cB)=s$. Then there exists a distribution whose matrix form $\cC$ on $M_{0..n}$ of seed length $s+O(\log(w/\delta))$ such that
$$\nmax{\inp{\cC}-\inp{\cA}\inp{\cB}}\le \delta.$$
\end{lemma}
This lemma is usually achieved with a pseudorandom object. For example, the INW generator~\cite{INW94} uses a bipartite expander with degree $\poly(w/\delta)$ to construct the distribution $\cC$ in the above lemma. That is, for every edge $(x,y)$ in the expander $G$, they  add $\cA(x)\cB(y)$ into $\cC$. A similar lemma can also be obtained with universal hash functions~\cite{Nis92} or seeded extractors~\cite{RR99}. By recursively constructing good approximations of $M_{0..n/2}$ and $M_{n/2..n}$ and applying \Cref{lemma:INW}, one can obtain a PRG which has seed length $O(\log n\cdot\log(nw/\eps))$ ($\delta$ is taken to be $\eps/n$ because of a union bound). Observe that in such constructions, one needs to pay $O(\log(1/\eps))$ (in seed length) per level of recursion. 

The crucial idea in \cite{BCG18} is to amortize this cost over all $\log n$ levels. What makes this possible is the following argument, which we will refer to as the \emph{sampler argument}. First we define the notion of an averaging sampler. 
\begin{definition}
A function $g:\bits{n}\times\bits{d}\to\bits{m}$ is an $(\eps,\delta)$-(averaging) sampler if for every function $f:\bits{m}\to{[0,1]}$, $$\pr[x\in\bits{n}]{\abs{\ex[s\in\bits{d}]{f(g(x,s))}-\ex[y\in\bits{m}]{f(y)}}\le\eps}\ge 1-\delta.$$
\end{definition}
\noindent
The crucial observation in \cite{BCG18} is that if one uses a sampler to prove \Cref{lemma:INW}, the error actually scales with the norm of one of the matrix forms.
\begin{lemma}[\cite{BCG18}]\label{lemma:samp-arg}
Consider a $(n,w)$-ROBP with matrix representation $[M]_{[n]}^{\bit}$. Let $\cA$ and $\cB$ be (pseudo)distributions in matrix form on $M_{0..n/2}$ and $M_{n/2..n}$ respectively. Let $n=\sout(\cA)$, $m=\sout(\cB)$. Suppose $\forall x\in\bits{n},\norm{\cA(x)}\le 1$ and $\forall y\in\bits{m},\norm{\cB(y)}\le 1$. Let $g:\bits{n}\times \bits{d}\to\bits{m}$ be a $(\eps,\delta)$ sampler. Then there exists a (pseudo)distribution $\cC$ such that 
$$\norm{\inp{C}-\inp{A}\inp{B}}\le O\left(w^2\left(\delta+\eps \ex[x]{\norm{\cA(x)}}\right)\right).$$
Besides, $\cC$ has outer seed length $n=\sout(\cA)$, and for every $x\in\bits{n}$, $$\cC(x)=\ex[s]{\cA(x)\cB\left(g(x,s)\right)}.$$
Note that $\si(\cC)=\si(\cA)+\si(\cB)+d$.
\end{lemma}
The intuition behind this approximation is as follows. If we want to compute the matrix product precisely, we take every $\cA(x)$ and multiply it with $\ex[y]{\cB(y)}$. However, with the help of sampler, we can use $x$ as our seed to select some samples from $\cB$, and take their average as an estimate of $\ex[y]{\cB(y)}$. The error of this approximation comes in two different way. For those $x$ which are not good choices of a seed for the sampler, the samples chosen with such an $x$ can deviate from the average arbitrarily. However, only $\delta$ fraction of $x$ can be bad, so they incur at most $\delta$ error. The second kind of error is the estimation error between average of samples $\ex[s]{\cB(g(x,s))}$ and the real average $\ex[y]{\cB(y)}$, which can be at most $\eps$. Since this gets multiplied  with $\cA(x)$, this kind of error actually scales with $\norm{\cA(x)}$. Although the first kind of error (which is $\delta$) does not benefit from $\norm{\cA}$ being small, in \cite{BCG18} they observe that, the parameter $\delta$ has almost no influence on the seed length in some cases. To discuss this more precisely, we first recall explicit constructions of samplers.
\begin{lemma}[\cite{RVW01,Gol11}]
For every $\delta,\eps>0$ and integer $m$, there exists a space efficient $(\eps,\delta)$-sampler $f:\bits{n}\times\bits{d}\to\bits{m}$ s.t. $d=O(\log\log(1/\delta)+\log(1/\eps))$ and $n=m+O(\log(1/\delta))+O(\log(1/\eps))$.
\end{lemma}
Note that in \Cref{lemma:samp-arg}, $s(\cC)=s(\cA)+d+\si(\cB)$. Therefore if $n\ge m+O(\log(1/\delta))+O(\log(1/\eps))$, $\delta$ has almost no impact on the seed length. 

To use the above ideas, it boils down to working with matrices with small norm, and  making sure that every multiplication is ``unbalanced" enough so that $\delta$ has no impact. \cite{BCG18} applies a delicate telescoping sum trick (which they called ``delta sampler") to divide an $\eps$-approximation into a base approximation with $1/\poly(n)$ error and several ``correcting terms" which have small norm. By carefully choosing the samplers and discarding all the non-necessary terms,  they roughly ensure the following properties: first, a matrix with large seed length must have small norm; second, every matrix multiplication is unbalanced enough so that $\delta$ has no impact on the seed length. 


With these properties and the sampler argument, they show that the total seed length is bounded by $\tilde{O}(\log(1/\eps)+\log n \log(nw))$.

\subsection{Our construction}\label{sec:our-const}
In executing the ideas sketched above,  the construction and analysis in \cite{BCG18} turns out to be quite complicated and involved. One thing which complicates the construction and analysis is its \emph{bottom-up} nature. That is, when multiplying two terms, they create more terms with the telescoping sum trick. Moreover, in the telescoping sum trick one needs to choose the parameters of each sampler very carefully to make sure the seed length of each term does not exceed its ``smallness". 

Our first step toward a simpler construction is the following \emph{top-down} formula, which we will apply recursively to compute an approximation of $M_{0..n}$:
\begin{lemma}\label{lemma:main}
Let $\norm{\cdot}$ be a sub-multiplicative matrix norm, and $A,B$ be two matrices s.t. $\norm{A},\norm{B}\le 1$. Let $k\in\bbN$ and $\gamma<1$. For every $0\le i\le k$, let $A_i$ be a $\gamma^{i+1}$-approximation of $A$, and let $B_i$ be a $\gamma^{i+1}$-approximation of $B$. Then $$\sum_{i=0}^k A_{i}B_{k-i}-\sum_{i=0}^{k-1} A_{i}B_{k-1-i}$$ is a $((k+2)\gamma^{k+1}+(k+1)\gamma^{k+2})$-approximation of $AB$.
\end{lemma}
\begin{proof}
We have,
\begin{align*}
&\norm{ (\sum_{i=0}^k A_{i}B_{k-i}-\sum_{i=0}^{k-1} A_{i}B_{k-1-i})-AB}\\
 &= \norm{ \sum_{i=0}^k (A-A_{i})(B-B_{k-i})-\sum_{i=0}^{k-1} (A-A_{i})(B-B_{k-1-i})+(A_k-A)B+A(B_k-B)}\\
 &\le \sum_{i=0}^k \norm{A-A_i}\cdot\norm{B-B_{k-i}}+\sum_{i=0}^{k-1}\norm{A-A_{i}}\cdot\norm{B-B_{k-1-i}}\\
 &+\norm{A_k-A}\cdot\norm{B}+\norm{A}\cdot\norm{B_k-B}\\
 &\le(k+2)\gamma^{k+1}+(k+1)\gamma^{k+2}
\end{align*}
\end{proof}
This formula shares an important property with the BCG construction: we never need a $\gamma^k$-approximation (which implies large seed length) on both sides simultaneously. The benefit of our top-down formula is that we are treating the PRPD as one object instead of the sum of many different terms. One obvious effect of such treatment is we don't need to analyze the ``smallness" of each term and the accuracy of the whole PRPD separately. 

In this top-down formula, we do not explicitly maintain small-norm matrices as in \cite{BCG18}. However, observe that in the proof of \Cref{lemma:main}, we are using the fact that $A_k-A$ is a small norm matrix. Our goal is to apply the sampler argument (\Cref{lemma:samp-arg}) on these ``implicit" small-norm matrices. The following is our main technical lemma.

\begin{lemma}[main lemma, informal]\label{lemma:main-informal}
Let $A,B\in\bbR^{w\times w}$, $k\in\bbN$ and $\gamma<1$. Suppose for every $i\le k$ there exists pseudodistribution $\cA_i,\cB_i$ such that $\ex[x]{\norm{\cA_i(x)-A}}\le \gamma^{i+1}$, $\ex[x]{\norm{\cB_i(x)-B}}\le \gamma^{i+1}$, and $\norm{\cA_i(x)},\norm{\cB_i(x)}\le 1$ for every $x$. Then there exists a pseudo-distribution $\cC_k$ such that $$\ex[x]{\norm{\cC_k(x)-AB}\le O(\gamma)^{k+1}},$$
where $\cC_k(x)=\sum_{i+j=k} A_{x,i}B_{x,j}-\sum_{i+j=k-1} A_{x,i}B_{x,j}$. $A_{x,i}$ and $B_{x,i}$ are defined as follows.
\begin{itemize}
\item 
If $i>\ceil{k/2}$, $A_{x,i}=\cA_i(x)$ and $B_{x,i}=\cB_i(x)$.
\item
If $i\le \ceil{k/2}$, $A_{x,i}=\ex[s]{\overline{\cA_i}(g_i(x,s))}$ and $B_{x,i}=\ex[s]{\overline{\cB_i}(g_i(x,s))}$, where $g_i$ is a $(\gamma^{i+1},\gamma^{k+1})$-sampler, and $\overline{\cA_i},\overline{\cB_i}$ denote the ``flattened" form of $\cA_i$ and $\cB_i$.
\end{itemize}
\end{lemma}

We leave the explanation of ``flattened" for later and explain the intuition behind the lemma first. Our goal is to construct $\cC_k$ such that $\cC_k(x)$ is a good approximation of $AB$ \emph{on average} over $x$. We  know that $\cA_i$ and $\cB_i$ are $\gamma^{i+1}$-approximation of $A$ and $B$ on average. Our hope is to use $x$ to draw samples $A_i$ and $B_i$ from $\cA_i$ and $\cB_i$, and apply the formula in \Cref{lemma:main} to get a good approximation of $AB$. In particular, a natural choice would be setting $A_{x,i}=\cA_i(x)$ and $B_{x,i}=\cB_i(x)$ for every $i\le k$. However, if there exists a term $A_{x,i}B_{x,j}$ such that $A_{x,i}$ and $B_{x,j}$ are both bad approximation for a large enough fraction of $x$, we cannot guarantee to get a $O(\gamma^{k+1})$-approximation on average.

To avoid the above case, for every $i\le\ceil{k/2}$ we use a sampler to approximate $\inp{\cA_i}$ and $\inp{\cB_i}$. This ensure that the chosen samples $A_{x,i}$ and $B_{x,i}$ are good \emph{with high probability}. This guarantees that in each term $A_{x,i}B_{x,j}$, at least one of $A_{x,i}$ or $B_{x,j}$ will be a good choice with high probability over $x$. If $A_{x,i}$ is a good choice with high probability, we can apply the average-case guarantee on $B_{x,i}$ to get an average-case guarantee for $\cC_k$, and vice versa. (Indeed, this is the sampler argument.) Therefore we can ensure that $\cC_k(x)$ is good on average. Note that we only apply a sampler on $\cA_i$ (or $\cB_i$) when $i\le \ceil{k/2}$, which means $\cA_i$ (or $\cB_i$) has small seed length. Therefore we don't need to add too much redundant seed to make the sampler argument work.

In executing the above sketched idea, we run into the following problem: in each multiplication, the inner seed on both sides  aggregates to the upper level. If we start with pseudodistributions with non-zero inner seed in the bottom level, the inner seed would become $\Omega(n)$ in the topmost level. Therefore we need a way to limit the aggregation of inner seed.

In \cite{BCG18}, they run into a similar problem. To deal with this,  they apply a different multiplication rule, ``outer product", in some special cases to deal with this. However, the outer product does not seem applicable in our construction. Nevertheless, we observe that whenever we use a sampler to select matrix $A_{x,i}$, we only care about whether $\inp{\cA_i}$ is close to $A$, and we don't need most of $\cA_i(x)$ to be close to $A$ anymore. Therefore we will ``flatten" $\cA_i$ whenever we apply a sampler. That is, recall that each $\cA_i(x)$ is realized by the average of some matrices, $\ex[y]{\cA_i(x,y)}$. We define the flattened form of $\cA_i$, denoted by $\overline{\cA}_i$, such that $\overline{\cA}_i(x\Vert y)=\cA_i(x,y)$. Observe that $\inp{\overline{\cA_i}}=\inp{\cA_i}$ and $\si(\overline{\cA_i})=0$. This guarantees that the inner seed length of $\cA_i$ will not aggregate in $\cC_k$. Moreover, while the flattening will increase the outer seed length of $\overline{\cA_i}$, this is almost for free since we only flatten $\cA_i$ when $i\le\ceil{k/2}$, i.e. when $\cA_i$ has relatively small seed length. As a result, this operation also helps us save a $O(\log\log(1/\eps))$ factor in the seed length.

We  conclude by briefly discussing the seed length analysis. First note that we set $\gamma=1/\poly(n)$ to make sure that the error is affordable after a union bound. Now consider the inner seed length. Consider a term $A_iB_j$ such that $i\ge j$. In this term, part of the inner seed of $\cC$ is passed to $\cA_i$, and the other is used for the sampler on $B_j$. Since the seed length of the sampler only needs to be as large as the ``precision gap" between $\cA_i$ and $\cC_k$, the inner seed length of $\cC_k$ can be maintained at roughly $O(k\log(1/\gamma))=O(\log(1/\eps))$. However, after each multiplication, there's actually a $O(\log (nw/\gamma))=O(\log (nw))$ additive overhead. Note that this is necessary since the $k=0$ case degenerates to the INW generator. Therefore after $\log n$ levels of recursion, the inner seed length will be $O(\log(1/\eps)+\log n\cdot \log (nw))$. 

Besides, we also need the outer seed length of $\cC_k$ to be long enough so that we can apply a sampler on $\cA_{\ceil{k/2}}$ and $\cB_{\ceil{k/2}}$. The seed length caused by approximation accuracy $\eps$ can be bounded similarly as the inner seed length. However, the $O(\log n\cdot \log (nw))$ inner seed length will be added to the outer seed length several times, because of the flattening operation. Nevertheless, since we only do flattening for $\cA_{i}$ and $\cB_{i}$ where $i\le \ceil{k/2}$, this ensures that the flattening operation happens at most $\log k$ times. So the total outer seed length will be bounded by $O(\log(1/\eps)+\log k\cdot \log n\cdot \log (nw))=O(\log(1/\eps)+\log \log(1/\eps)\cdot \log n\cdot \log (nw))$, which is bounded by $O(\log(1/\eps)+\log \log(nw)\cdot \log n\cdot \log (nw))$ since $O(\log(1/\eps))$ is the dominating term when $\log(1/\eps)\ge\log^3(nw)$.


\section{Preliminaries}\label{sec:prelim}
\subsection{Averaging samplers}
\begin{definition}
A function $g:\bits{n}\times\bits{d}\to\bits{m}$ is a \emph{$(\eps,\delta)$ (averaging) sampler} if for every function $f:\bits{m}\to{[0,1]}$, 
$$\pr[x\in\bits{n}]{\abs{\ex[s\in\bits{d}]{f(g(x,s))}-\ex[y\in\bits{m}]{f(y)}}\le\eps}\ge 1-\delta.$$
\end{definition}
\noindent
It's easy to show that samplers also work for $f$ with general range by scaling and shifting.
\begin{claim}\label{lemma:sampler-general}
Let $g:\bits{n}\times\bits{d}\to\bits{m}$ be a $(\eps,\delta)$-sampler, and let $\ell<r\in\bbR$. Then for every $f:\bits{m}\to{[\ell,r]}$, 
$$\pr[x\in\bits{n}]{\abs{\ex[s\in\bits{d}]{f(g(x,s))}-\ex[y\in\bits{m}]{f(y)}}\le\eps(r-\ell)}\ge 1-\delta.$$
\end{claim}
\begin{proof}
Let $f'$ be the function such that $f'(y)=(f(y)-\ell)/(r-\ell)$. Observe that the range of $f'$ is in $[0,1]$. By definition of sampler, 
$$\pr[x\in\bits{n}]{\abs{\ex[s\in\bits{d}]{f'(g(x,s))}-\ex[y\in\bits{m}]{f'(y)}}\le\eps}\ge 1-\delta.$$
By multiplying $(r-\ell)$ on both sides of the inequality inside the probability above we prove the claim.
\end{proof}
In our construction, we will use the following sampler which is explicitly computable with small space.
\begin{lemma}[\cite{RVW01,Gol11}]\label{lemma:samp-opt}
For every $\delta,\eps>0$ and integer $m$, there exists a $(\eps,\delta)$-sampler $f:\bits{n}\times\bits{d}\to\bits{m}$ s.t. $d=O(\log\log(1/\delta)+\log(1/\eps))$ and $n=m+O(\log(1/\delta))+O(\log(1/\eps))$. Moreover, for every $x,y$, $f(x,y)$ can be computed in space $O(m+\log(1/\delta)+\log(1/\eps))$.
\end{lemma}
\begin{remark}
The original sampler in \cite{RVW01} has a restriction on $\eps$. Such a restriction will cause a $2^{O(\log^*(nw/\eps))}$ factor in our construction, as in \cite{BCG18}. However, \cite{RVW01} pointed out that the restriction is inherited from the extractor in \cite{Zuc97}, which breaks down when the error is extremely small. As observed in \cite{Gol11}, this restriction can be removed by plugging in a more recent extractor construction in \cite{GUV09}. Note that there exists a space-efficient implementation of \cite{GUV09} in \cite{KNW08}, so the resulting sampler is also space-efficient. For completeness we include a proof in \Cref{sec:proof-of-samp}.
\end{remark}
\subsection{Matrix norms}
As in \cite{BCG18}, we will use the infinity norm in this paper.
\begin{definition}
For every matrix $A\in\bbR^{w\times w}$, $\ninf{A}=\max_i{\sum_j |A_{i,j}|}$.
\end{definition}
We record some well known properties of the infinity norm.
\begin{claim}
Let $A,B\in\bbR^{w\times w}$, $c\in \bbR$. Then
\begin{itemize}
\item
$\norm{cA}=|c|\norm{A}$
\item
$\norm{A}+\norm{B}\le\norm{A+B}$
\item
$\ninf{AB}\le \ninf{A}\ninf{B}$ 
\item
 $\max_{i,j} |A_{i,j}|\le\ninf{A}$
\item
If $A$ is stochastic, then $\ninf{A}=1$
\end{itemize}
\end{claim}
Note that for any $(n,w)$-ROBP represented by $w\times w$ matrices $M_{[n]}^{\bit}$, $\ninf{M_{i..j}}=1$ for every $0\le i\le j\le n$.


\section{Approximate Matrix Multiplication via Samplers}\label{sec:samp-our}




In this section we formally prove the sampler arguments which will be used in our construction. Our proof strategy resembles that of \cite{BCG18},  with the following two crucial differences. First, we will define two different notions of ``smallness" for our flattening idea. Second, in our construction we need the case where we use samplers to select matrices on both sides (\Cref{lemma:sym}). 

We will consider mappings $\cA:\bits{n}\to\bbR^{w\times w}$ which correspond to the implicit small norm matrices we discussed in the previous section. Borrowing notation from \Cref{def:MBS}, we use $\inp{\cA}$ to denote $\ex[x]{\cA(x)}$. First we define two different norms for the mapping $\cA$. The \emph{robust norm} is similar to the notion of ``smallness" in \cite{BCG18}, i.e. the average of norm of $\cA(x)$, while the \emph{norm} of $\cA$ is simply the norm of $\inp{\cA}$, i.e. the norm of average of $\cA(x)$. 
\begin{definition}
For every function $\cA:\bits{n}\to\bbR^{w\times w}$, we define the \emph{norm} of $\cA$ to be $\norm{\cA}=\ninf{\ex[x\in\bits{n}]{\cA(x)}}$, and the \emph{robust norm} of $\cA$ to be  $\norm[r]{\cA}=\ex[x\in\bits{n}]{\ninf{\cA(x)}}$. Besides, we define the weight of $\cA$ to be  $\wt(\cA)=\max_x \ninf{\cA(x)}$.
\end{definition}
\begin{claim}\label{claim:norms}
$\norm{\cA}\le\norm[r]{\cA}\le \wt(\cA)$.
\end{claim}
\begin{proof}
$\norm{\cA}\le\norm[r]{\cA}$ is by sub-additivity of $\norm{\cdot}$, and $\norm[r]{\cA}\le \wt(\cA)$ since $\norm[r]{\cA}$ is the average of values no larger than $\wt(\cA)$.
\end{proof}
\noindent
Next we show a simple lemma which will be used later. That is, a sampler for functions with range $[0,1]$ is also a sampler for matrix-valued functions, where the error is measured with infinity norm. 
\begin{lemma}\label{lemma:mat-samp-general}
For every function $\cA:\bits{m}\to\bbR^{w\times w}$ and every $(\eps,\delta)$-sampler $g:\bits{n}\times\bits{d}\to\bits{m}$,
$$\pr[x\in\bits{n}]{\ninf{\ex[s\in\bits{d}]{\cA(g(x,s))}-\inp{\cA}}\le 2w\wt(\cA)\eps}\ge 1- w^2\delta.$$
\end{lemma}
\begin{proof}
Let $\cE(y)=\cA(y)-\inp{\cA}$. For every $i,j\in[w]$, observe that 
$$\max_y \cE(y)_{i,j}-\min_y \cE(y)_{i,j}= \max_y \cA(y)_{i,j}-\min_y \cA(y)_{i,j}$$ 
By the property of sampler it follows that
$$\pr[x\in\bits{n}]{\abs{\ex[s]{\cE(g(x,s))_{i,j}}}\le 2\eps\wt(\cA)}\ge 1-\delta.$$ Using a union bound, 
$$\pr[x\in\bits{n}]{\forall i,j\in[w], \abs{\ex[s]{\cE(g(x,s))_{i,j}}}\le 2\eps\wt(\cA) }\ge 1- w^2\delta.$$
Thus by definition of the infinity norm, we can conclude that 
$$\pr[x\in\bits{n}]{\ninf{\ex[s\in\bits{d}]{\cE(g(x,s))}}\le 2w\wt(\cA)\eps }\ge 1-w^2\delta.$$
which by sub-additivity of $\ninf{\cdot}$ implies 
$$\pr[x\in\bits{n}]{\ninf{\ex[s]{\cA(g(x,s))}}\le \norm{\cA}+2w\wt(\cA)\eps}\ge 1- w^2\delta.$$
\end{proof}
\begin{corollary}\label{lemma:mat-samp}
For every function $\cA:\bits{m}\to\bbR^{w\times w}$ and every $(\eps,\delta)$-sampler $g:\bits{n}\times\bits{d}\to\bits{m}$,
$$\pr[x\in\bits{n}]{\ninf{\ex[s\in\bits{d}]{\cA(g(x,s))}}\le \norm{\cA}+2w\wt(\cA)\eps}\ge 1- w^2\delta.$$
\end{corollary}
\begin{proof}
By sub-additivity of $\ninf{\cdot}$, $\ninf{\ex[s\in\bits{d}]{\cA(g(x,s))}-\inp{A}}\le 2w\wt(\cA)\eps$ implies $\ninf{\ex[s\in\bits{d}]{\cA(g(x,s))}}\le \norm{\inp{\cA}}+2w\wt(\cA)\eps$. The claim now directly follows from \Cref{lemma:mat-samp-general}.
\end{proof}
\noindent
Now we introduce three different matrix multiplication rules. The first one is applying a sampler on both sides, and the second and third are applying sampler on only one side. 
\begin{lemma}[symmetric product]\label{lemma:sym}
Consider $\cA:\bits{n}\to\bbR^{w\times w}$ and $\cB:\bits{m}\to\bbR^{w\times w}$. Let $f:\bits{k}\times \bits{d_A}\to\bits{n}$ be a $(\delta,\eps_A)$ sampler, and $g:\bits{k}\times \bits{d_B}\to\bits{m}$ be a $(\delta,\eps_B)$ sampler. Then 
$$\ex[z]{\ninf{\ex[x,y]{\cA(f(z,x))\cB(g(z,y))}}}\le 2w^2\delta \wt(\cA)\wt(\cB)+\left(\norm{\cA}+2w\wt(\cA)\eps_A\right)\left(\norm{\cB}+2w\wt(\cB)\eps_B\right).$$
\end{lemma}
\begin{proof}
Let $$E_A=\left\{z: \ninf{\ex[x]{\cA(f(z,x)}}> \norm{\cA}+2w\wt(\cA)\eps_A\right\},$$and $$E_B=\left\{z: \ninf{\ex[y]{\cB(g(z,y)}}> \norm{\cB}+2w\wt(\cB)\eps_B\right\}.$$
Define $E=E_A\cup E_B$. By Lemma~\ref{lemma:mat-samp} and union bound, $\pr[z]{z\in E}<2w^2\delta$. Therefore
\begin{align*}
\ex[z]{\ninf{\ex[x,y]{\cA(f(z,x))\cB(g(z,y))}}}
&=\pr{z\in E}\ex[z\in E]{\ninf{\ex[x,y]{\cA(f(z,x))\cB(g(z,y))}}} \\
&+ \pr{z\not\in E}\ex[z\not\in E]{\ninf{\ex[x]{\cA(f(z,x))}\ex[y]{\cB(g(z,y))}}}\\
&\le 2w^2\delta \wt(\cA)\wt(\cB)+\ex[z\not\in E]{\ninf{\ex[x]{\cA(f(z,x))}}\ninf{\ex[y]{\cB(g(z,y))}}}\\
&\le 2w^2\delta \wt(\cA)\wt(\cB)+\left(\norm{\cA}+2w\wt(\cA)\eps_A\right)\left(\norm{\cB}+2w\wt(\cB)\eps_B\right).
\end{align*}
The second last inequality is by the fact that $\norm{\cdot}$ is non-negative and sub-multiplicative.
\end{proof}
\begin{lemma}[left product]\label{lemma:left}
Consider $\cA:\bits{k}\to\bbR^{w\times w}$ and $\cB:\bits{m}\to\bbR^{w\times w}$. Let $g:\bits{k}\times \bits{d_B}\to\bits{m}$ be a $(\delta,\eps_B)$ sampler. Then 
$$\ex[z]{\ninf{\ex[y]{\cA(z)\cB(g(z,y))}}}\le w^2\delta \wt(\cA)\wt(\cB)+\norm[r]{\cA}\left(\norm{\cB}+2w\wt(\cB)\eps_B\right).$$
\end{lemma}
\begin{proof}
Let $$E=\left\{z: \ninf{\ex[y]{\cB(g(z,y)}}> \norm{\cB}+2w\wt(\cB)\eps_B\right\}.$$
By Lemma~\ref{lemma:mat-samp}, $\pr[z]{z\in E}<w^2\delta$. Therefore 
\begin{align*}
\ex[z]{\ninf{\ex[y]{\cA(z)\cB(g(z,y))}}}
&=\pr{z\in E}\ex[z\in E]{\ninf{\ex[y]{\cA(z)\cB(g(z,y))}}} \\
&+ \pr{z\not\in E}\ex[z\not\in E]{\ninf{\ex[y]{\cA(z)\cB(g(z,y))}}}\\
&\le w^2\delta \wt(\cA)\wt(\cB)+\pr{z\not\in E}\cdot\ex[z\not\in E]{\ninf{\cA(z)}\ninf{\ex[y]{\cB(g(z,y))}}}\\
&\le w^2\delta \wt(\cA)\wt(\cB)+\pr{z\not\in E}\ex[z\not\in E]{\ninf{\cA(z)}}\cdot(\norm{\cB}+2w\wt(\cB)\eps_B)\\
&\le w^2\delta \wt(\cA)\wt(\cB)+\norm[r]{\cA}\left(\norm{\cB}+2w\wt(\cB)\eps_B\right).
\end{align*}
The third last inequality is by sub-multiplicativity of $\norm{\cdot}$, the second last inequality is by non-negativity of $\norm{\cdot}$, and the last inequality is by the fact that $$\pr{z\not\in E}\cdot\ex[z\not\in E]{\norm{\cA(z)}}=\ex[z]{\norm{\cA(z)}\cdot \mathbbm{1}(z\not\in E)}\le \norm[r]{\cA}.$$
\end{proof}
\begin{lemma}[right product]\label{lemma:right}
Consider $\cA:\bits{k}\to\bbR^{w\times w}$ and $\cB:\bits{m}\to\bbR^{w\times w}$. Let $f:\bits{k}\times \bits{d_A}\to\bits{n}$ be a $(\delta,\eps_A)$ sampler. Then 
$$\ex[z]{\ninf{\ex[x]{\cA(f(z,x))\cB(z)}}}\le w^2\delta \wt(\cA)\wt(\cB)+\left(\norm{\cA}+2w\wt(\cA)\eps_A\right)\norm[r]{\cB}.$$
\end{lemma}
\begin{proof}
Let $$E=\left\{z: \ninf{\ex[x]{\cA(f(z,x)}}> \norm{\cA}+2w\wt(\cA)\eps_A\right\}.$$
By Lemma~\ref{lemma:mat-samp}, $\pr[z]{z\in E}<w^2\delta$. Therefore 
\begin{align*}
\ex[z]{\ninf{\ex[x]{\cA(f(z,x))\cB(z)}}}
&=\pr{z\in E}\ex[z\in E]{\ninf{\ex[x]{\cA(f(z,x))\cB(z)}}} \\
&+ \pr{z\not\in E}\ex[z\not\in E]{\ninf{\ex[x]{\cA(f(z,x))\cB(z)}}}\\
&\le w^2\delta \wt(\cA)\wt(\cB)+\pr{z\not\in E}\cdot\ex[z\not\in E]{ \ninf{\ex[x]{\cA(f(z,x))}}\ninf{\cB(z)}}\\
&\le w^2\delta \wt(\cA)\wt(\cB)+(\norm{\cA}+2w\wt(\cA)\eps_A)\cdot \pr{z\not\in E}\ex[z\not\in E]{\ninf{\cB(z)}}\\
&\le w^2\delta \wt(\cA)\wt(\cB)+\left(\norm{\cA}+2w\wt(\cA)\eps_A\right)\norm[r]{\cB}.
\end{align*}
\end{proof}
\section{Main Construction}\label{sec:main-construction}
In this section we show our main construction and prove its correctness. We first introduce several definitions.
\begin{definition}
For every mapping $\cA:\bits{n}\to\bbR^{w\times w}$ and every matrix $A\in\bbR^{w\times w}$, we define $\cA-A$ to be the mapping s.t. $(\cA-A)(x)=\cA(x)-A$.
\end{definition}
\begin{definition}
Consider $A\in\bbR^{w\times w}$ and $\cA:\bits{n}\to\bbR^{w\times w}$. $\cA$ is a $\eps$-\emph{approximator} of $A$ if $\norm{\ex[x]{\cA(x)}-A}\le \eps$, i.e. $\norm{\cA-A}\le \eps$. $\cA$ is a $\eps$-\emph{robust approximator} of $A$ if $\ex[x]{\norm{\cA(x)-A}}\le \eps$, i.e. $\norm[r]{\cA-A}\le \eps$.
\end{definition}
\noindent
Now we define a robust PRPD. Note that a $(n,w,\eps)$-robust PRPD $(G,\rho)$ is also a $\wt(G,\rho)$-bounded $(n,w,\eps)$-PRPD.
\begin{definition}\label{def:robust-PRPD}
$(G,\rho):\bits{\sout}\times\bits{\si}\times{[\wt]}\to\bits{n}\times\bbR$ is a $(n,w,\eps)$-\emph{robust PRPD} if for every $(n,w)$-ROBP and its matrix representation $[M]_{[n]}^{\bit}$ the following holds. Let $\cA:\bits{\sout}\times\bits{\si}\to\bbR^{w\times w}$ denote the mapping 
$$\cA(x,y)=\ex[i\in {[\wt]}]{\rho(x,y,i)\cdot M_{0..n}^{G(x,y,i)}}.$$
\begin{itemize}
\item
Every $\rho(x,y,i)$ is either $\wt$ or $-\wt$. In other word, $\cA(x,y)$ is the summation of transition matrices with coefficient $\pm 1$.
\item
Let $\widehat{\cA}$ denote the mapping $\widehat{\cA}(x)=\ex[y]{\cA(x,y)}$. Then $\widehat{\cA}$ is a $\eps$-robust approximator for $M_{0..n}$. 
\end{itemize}
We say $\wt$ is the weight of $(G,\rho)$, denoted by $\wt(G,\rho)$. $\sout$ is the \emph{outer seed length} of $(G,\rho)$, denoted by $\sout(G,\rho)$. $\si$ is the \emph{inner seed length} of $(G,\rho)$, denoted by $\si(G,\rho)$. We write $s(G,\rho)=\sout(G,\rho)+\si(G,\rho)$ for short. We say $(G,\rho)$ is explicit if it can be computed in $O(s(G,\rho))$ space.

We say $\cA$ is the \emph{matrix form} of $(G,\rho)$ on $M_{0..n}$, and the definition of $\sout,\si,\wt$ on $(G,\rho)$ also apply to $\cA$. We say $\widehat{\cA}$ is the \emph{robust matrix form} of $(G,\rho)$ on $M_{0..n}$. 
\end{definition}
\begin{remark}\label{rmk:mon}
The above definition is similar to \Cref{def:MBS}, but each matrix $\cA(x,y)$ is realized with $\wt$ matrices instead of one matrix. These $\wt$ matrices will never be separated even after flattening. We do this in order to ensure that the matrix form always take bit-strings as input. This ensures that we can increase the outer and inner seed length of $\cA$ arbitrarily: we can construct the new mapping $\cA':\bits{\sout'}\times\bits{\si'}$ such that $\cA'(x,y)=\cA(x_p,y_p)$ where $x_p$ is the length-$\sout(\cA)$ prefix of $x$ and $y_p$ is the length-$\si(\cA)$ prefix of $y$. In other word, $\cA'$ computes the output only with prefix of necessary length of the input, and ignore the remaining bits. It is easy to verify that $\cA'$ is also the matrix form of a $(n,w,\eps)$-robust PRPD.
\end{remark}
\noindent
The following is some additional basic properties about robust PRPD and its flattened form.
\begin{claim}\label{claim:std-approx}
Let $(G,\rho):\bits{\sout}\times\bits{\si}\times{[\wt]}\to\bits{n}\times\bbR$ be a $(n,w,\eps)$-\emph{robust PRPD}. For every $(n,w)$-ROBP $M_1^0,M_1^1,\dots,M_n^0,M_n^1$ the following holds.
\begin{itemize}
\item 
Let $\widehat{\cA}$ be the robust matrix form of $(G,\rho)$ on $M_{0..n}$. Then $\wt(\widehat{\cA})\le \wt(G,\rho)$.
\item
Let $\cA$ denote the matrix form of $(G,\rho)$ on $M_{0..n}$. Let $\overline{\cA}:\bits{\sout+\si}\to\bbR^{w\times w}$ denote the mapping $\overline{\cA}(x\Vert y)=\cA(x,y)$. We say $\overline{\cA}$ is the \emph{flattened matrix form} of $(G,\rho)$ on $M_{0..n}$. Then $\overline{\cA}$ is an $\eps$-approximator for $M_{0..n}$, and $\wt(\overline{\cA})\le\wt(G,\rho)$.
\end{itemize}
\end{claim}
\begin{proof}
Recall that for every string $r\in\bits{n}$, $\norm{M_{0..n}^r}=1$. By sub-additivity of $\norm{\cdot}$ we have $\norm{\cA(x,y)}\le\wt(G,\rho)$ for every $x,y$, which implies $\wt(\overline{\cA})\le\wt(G,\rho)$. By sub-additivity and scalibility of $\norm{\cdot}$, we have $\wt(\cA')\le \wt(\cA)$. To show that $\overline{\cA}$ is a $\eps$-approimxator of $M_{0..n}$, observe that $\cA'$ is also an $\eps$-approximator of $M_{0..n}$ by \Cref{claim:norms}, and note that $\inp{\cA}=\inp{\cA'}$.
\end{proof}
\noindent
Now we prove our main lemma. The following lemma allows us to construct robust PRPDs for $(2m,w)$ ROBPs from robust PRPDs for $(m,w)$ ROBPs, without increasing the seed length too much. We will recursively apply this lemma for $\log n$ levels to get a $(n,w,\eps)$-robust PRPD. The basic idea is as described in \Cref{lemma:main-informal}.
\begin{lemma}\label{lemma:main-approx}
Suppose there exists $\sout,\si$ such that the following conditions hold.
\begin{itemize}
\item
For every $0\le i\le k$, there exists a $(m,w,\gamma^{i+1})$-robust PRPD $(G_i,\rho_i)$ s.t. $\wt(G_i,\rho_i)\le \binom{m-1}{i}$ and $\sout(G,\rho)\le \sout$. Moreover, for every $0\le i\le \ceil{k/2}$, $s(G_i,\rho_i) \le \sout$.
\item
For every $i\le \ceil{k/2}$, there exists a $(\eps_i,\delta)$-sampler $g_i:\bits{\sout}\times\bits{d_i}\to\bits{s(G_i,\rho_i)}$, where $\eps_i\le \gamma^{i+1}/(w\cdot\binom{m-1}{i})$ and $\delta\le \gamma^{k+1}/(w^2\cdot \binom{2m-1}{i})$.
\item
For every $i\ge j\ge 0$ s.t. $i+j\le k$, if $j\le i\le \ceil{k/2}$, then $d_i+d_j\le \si$. If $i>\ceil{k/2}$, then $\si(G_i,\rho_i)+d_j\le \si $.
\end{itemize}
Then there exists a $(2m,w,(11\gamma)^{k+1})$-robust PRPD $(G,\rho)$ s.t. $\sout(G,\rho)=\sout$, $\si(G,\rho)=\si$ and $\wt(G,\rho)\le \binom{2m-1}{k}$.
\end{lemma}
\begin{proof}
Fix any $(2m,w)$-ROBP with matrix representation $M_{[2m]}^{\bit}$. Let $A=M_{0..m}$ and $B=M_{m..2m}$.
 For every $0\le i\le k$, let $\cA_i,\widehat{\cA}_i,\overline{\cA}_i$ denote the matrix form, robust matrix form and flattened matrix form of $(G,\rho)$ on $M_{0..m}$ respectively. Let $\cB_i,\widehat{\cB}_i,\overline{\cB}_i$ denote the matrix form, robust matrix form and flattened matrix form of $(G,\rho)$ on $M_{m..2m}$ respectively. By definition, $\widehat{\cA}_i$ and $\widehat{\cB}_i$ are $\gamma^{i+1}$-robust approximator for $A$ and $B$ respectively. By \Cref{claim:std-approx}, $\overline{\cA}_i$ and $\overline{\cB}_i$ are $\gamma^{i+1}$-approximator for $A$ and $B$ respectively. Moreover, we will increase the outer seed length of $\cA_i$ and $\cB_i$ to match the length of the given input when necessary. (See \Cref{rmk:mon})
 
Now for every $x,y$ we define a mapping $\cC_k:\bits{\sout}\times\bits{\si}\to\bbR^{w\times w}$ as follows. Note that $\cC_k$ corresponds to the matrix form of $(G,\rho)$ on $M_{0..2m}$.
\begin{enumerate}[(1)]
\item
For every $0\le i\le \ceil{k/2}$, let $a_i$ be the prefix of $y$ of length $d_i$ and $b_i$ be the suffix of $y$ of length $d_i$. Define $A_{x,y,i}=\overline{\cA_i}(g_i(x,a_i))$ and $B_{x,y,i}=\overline{\cB_i}(g_i(x,b_i))$.
\item
For every $\ceil{k/2}<i\le k$, let $a_i$ be the prefix of $y$ of length $\si(\cA_i)$ and $b_i$ be the suffix of $y$ of length $\si(\cB_i)$. Define $A_{x,y,i}=\cA_i(x,a_i)$ and $B_{x,y,i}=\cB_i(x,b_i)$.
\item
Define $\cC_k(x,y)=\sum_{i+j=k} A_{x,y,i}B_{x,y,j}-\sum_{i+j=k-1}A_{x,y,i}B_{x,y,j}$.
\end{enumerate}
Note that for every $i+j\le k$, prefix $a_i$ and suffix $b_j$ of $y$ never overlap.

By expanding every $A_{x,y,i}B_{x,y,j}$ term with distributive law, we can see that each small term in $A_{x,y,i}B_{x,y,j}$ has coefficient $\pm 1$, which satisfies the first condition of robust PRPD. Moreover, the total number of terms after expanding is 
$$\wt(\cC_k)\le \sum_{i+j=k}\binom{m-1}{i}\cdot \binom{m-1}{j} + \sum_{i+j=k-1}\binom{m-1}{i}\cdot \binom{m-1}{j}= \binom{2m-1}{k}.$$
It remains to show that $\cC_k$ satisfies the second condition of robust PRPD, i.e. $\ex[y]{\cC_k(x,y)}$ is a good approximation of $M_{0..2m}=AB$ on average over $x$. Observe that 
\begin{align*}
\ex[x]{\norm{\ex[y]{\cC_k(x,y)}-AB}}
&=\ex[x]{\norm{\ex[y]{\cC_k(x,y)-AB}}}\\
&\le \sum_{i+j=k} \ex[x]{\norm{\ex[y]{(A_{x,y,i}-A)(B_{x,y,j}-B)}}} \\
&+ \sum_{i+j=k-1} \ex[x]{\norm{\ex[y]{(A_{x,y,i}-A)(B_{x,y,j}-B)}}}\\
&+ \ex[x]{\norm{\ex[y]{(A_{x,y,k}-A)B}}} + \ex[x]{\norm{\ex[y]{A(B_{x,y,k}-B)}}},
\end{align*}
by decomposing $\cC_k(x,y)-AB$ with the equation in the proof of \Cref{lemma:main} and applying sub-additivity of $\norm{\cdot}$. 

First we consider the last two terms. Since $\norm{B}=1$, by sub-multiplicativity we have 
$$\ex[x]{\norm{\ex[y]{(A_{x,y,k}-A)B}}}\le \ex[x]{\norm{\ex[y]{A_{x,y,k}-A}}}.$$
Now consider two cases. If $k\ge 2$, then 
$$\ex[x]{\norm{\ex[y]{A_{x,y,k}-A}}}=\ex[x]{\norm{\widehat{\cA}_k(x)-A}}\le \gamma^{k+1}$$ by definition. If $k<2$, then 
$$\ex[x]{\norm{\ex[y]{A_{x,y,k}-A}\cdot B}}=\ex[x]{\norm{\ex[a_k]{\overline{\cA_k}(g_k(x,a_k))-A}}\cdot B}.$$
Apply \Cref{lemma:right} on $\overline{\cA_k}-A$ and the dummy mapping $\cB$ s.t. $\cB(x)=B$ for every $x$, we can derive that the above formula is bounded by $w^2\delta\binom{m-1}{k}+3\gamma^{i+1}$. For the term $\ex[x]{\norm{\ex[y]{A(B_{x,y,k}-B)}}}$ we can get the same bound with a similar proof.

Now consider the terms in the form $\ex[x]{\norm{\ex[y]{(A_{x,y,i}-A)(B_{x,y,j}-B)}}}$. \\
First consider the case $i,j\le\ceil{k/2}$. Then
\begin{align*}
&\ex[x]{\norm{\ex[y]{(A_{x,y,i}-A)(B_{x,y,j}-B)}}}\\
&=\ex[x]{\norm{\ex[a_i]{\overline{\cA_i}(g_i(x,a_i))-A}\ex[b_j]{\overline{\cB_k}(g_j(x,b_j))-B}}}\textrm{ (since }a_i,b_j\textrm{ don't overlap)}\\
&\le 2w^2\delta \cdot \binom{m-1}{i}\cdot\binom{m-1}{j}+9\gamma^{i+j+2}.\textrm{ (by \Cref{lemma:sym})}
\end{align*}
Next consider the case $i>\ceil{k/2},j\le \ceil{k/2}$. Then 
\begin{align*}
&\ex[x]{\norm{\ex[y]{(A_{x,y,i}-A)(B_{x,y,j}-B)}}}\\
&=\ex[x]{\norm{\widehat{\cA}_i(x)\cdot\ex[b_j]{\overline{\cB_k}(g_j(x,b_j))-B}}}\textrm{ (since }a_i,b_j\textrm{ don't overlap)}\\
&\le w^2\delta \cdot \binom{m-1}{i}\cdot\binom{m-1}{j}+3\gamma^{i+j+2}.\textrm{ (by \Cref{lemma:right})}
\end{align*}
Similarly for the case that $i\le \ceil{k/2},j>\ceil{k/2}$ we can show that 
$$\ex[x]{\norm{\ex[y]{(A_{x,y,i}-A)(B_{x,y,j}-B)}}}\le w^2\delta \cdot \binom{m-1}{i}\cdot\binom{m-1}{j}+3\gamma^{i+j+2}$$
by \Cref{lemma:left}. Finally, note that the case $i,j>\ceil{k/2}$ does not exist because $i+j\le k$. 

Taking the summation of all the cases, we get 
\begin{align*}
&\ex[x]{\norm{\ex[y]{\cC_k(x,y)}-AB}}\\
&\le 2w^2\delta\cdot\left(\sum_{i+j=k}\binom{m-1}{i}\binom{m-1}{j}+\sum_{i+j=k-1}\binom{m-1}{i}\binom{m-1}{j}+\binom{m-1}{k}\right)\\
&+ (k+1)\cdot 9\gamma^{k+2}+ k\cdot 9\gamma^{k+1} + 2\cdot 3\gamma^{k+1} \\
&\le 4w^2\delta \cdot \binom{2m-1}{k} + (9k+9)\gamma^{k+2} + (9k+6)\gamma^{k+1}\\
&\le (10k+11)\gamma^{k+1}\\
&\le (11\gamma)^{k+1}.
\end{align*}
Moreover, note that $AB=M_{0..2m}$, and the construction of $\cC_k$ does not depend on the matrices $M_{[2m]}^{\bit}$. (See \Cref{sec:matrices} for how the arithmetic operations in $\cC_k(x,y)$ are translated back to operations on pseudo-distributions.) Therefore there exists a $(2m,w,(11\gamma)^{k+1})$-robust PRPD $(G,\rho)$. 
\end{proof}
\noindent
Finally we analyze the seed length of the recursive construction, and present the main theorem.
\begin{theorem}
There exists an explicit $(n,w,\eps)$-robust PRPD $(G,\rho)$ such that 
\begin{itemize}
\item 
$\sout(G,\rho)= O\left(\log(1/\eps)+ \log n \log (nw)\log\left(\frac{\log(1/\eps)}{\log n}\right)\right)$
\item
$\si(G,\rho)= O\left(\log(1/\eps)+ \log n \log (nw)\log\left(\frac{\log(1/\eps)}{\log n}\right)\right)$
\item
$\wt(G,\rho)=\poly(1/\eps)$
\end{itemize}
Moreover, for every $B$ the approximator $\cG$ has the same corresponding pseudodistribution.
\end{theorem}
\begin{proof}
Let $c$ be the constant such that for every $\eps,\delta>0$ there exists a $(\eps,\delta)$-sampler $g:\bits{n}\times\bits{d}\to\bits{m}$ such that $n=m+c\log(1/\eps)+c\log(1/\delta)$ and $d=c\log (1/\eps)+c\log\log(1/\delta)$, as guaranteed in \Cref{lemma:samp-opt}. WLOG assume that $n$ is a power of $2$. Define $\gamma=1/n^{4}$. For every $0\le h \le \log n$, every $k\ge 0$, we will inductively prove that there exists a $(2^h,w,(11^h\gamma)^{k+1})$-robust PRPD $(G_{h,k},\rho_{h,k})$ with the following parameters.
\begin{itemize}
\item
If $k\le 1$, $\sout(G_{h,k},\rho_{h,k})\le h \cdot (3ck\log(n/\gamma) + 7c\log(w/\gamma))$
\item 
If $k>1$, $\sout(G_{h,k},\rho_{h,k})\le 4ck\log (n/\gamma)+(\ceil{\log k}+1)\cdot h \cdot (10c\log(w/\gamma))$
\item
If $k\le 1$, $\si(G_{h,k},\rho_{h,k})\le ck\log(n/\gamma) + 4c\log(w/\gamma)$
\item
If $k>1$, $\si(G_{h,k},\rho_{h,k})\le ck\log (n/\gamma)+ h \cdot (4c\log(kw/\gamma))$
\item
$\wt(G_{h,k},\rho_{h,k})\le \max(1,\binom{2^h-1}{k})$
\end{itemize}
We will write $\sout[,h,k]=\sout(G_{h,k},\rho_{h,k})$ and $\si[,h,k]=\sout(G_{h,k},\rho_{h,k})$ for short. First consider the terminal case $2k\ge 2^h$ or $h=0$. In this case we simply take $\sout[,h,k]=0$, $\si[,h,k]=2^h \le 2k$ and $\wt(G_{h,k},\rho_{h,k})=1$ s.t. $G_{h,k}(x,y,i)=y$ and $\rho_{h,k}(x,y,i)=1$. For the other cases, we show that we can get the intended parameters by constructing $(G_{h,k},\rho_{h,k})$ with the recursion in \Cref{lemma:main-approx}. Note that based on the induction hypothesis we can assume $\cG_{a,h-1,k}$ and $\cG_{a+2^{h-1},h-1,k}$ have exactly the same parameters, so we consider the parameter of $\cG_{a,h-1,k}$ only. We have seen that the bound for $\wt(G_{h,k},\rho_{h,k})$ is correct. First we show that the bound for $\si[,h,k]$ is correct. Recall that in the recursion we take parameters $d_i=c\log(1/\eps_i)+c\log\log(1/\delta)\le ci\log(n/\gamma)+2c\log(knw/\gamma)$, based on the fact that $\binom{2^h-1}{i}\le n^i$. Now consider the restriction on $\si(\cG_k)$ in our recursion. For $i+j\le k$ and $j\le i\le\ceil{k/2}$, we need 
$$d_i+d_j\le ck\log(n/\gamma)+4c\log(knw/\gamma)\le\si[,h,k]$$ which is true. For $i+j\le k$ and $i>\ceil{k/2}$, we need 
\begin{align*}
\si[,h-1,i]+d_j
&\le ci\log(1/\gamma)+(h-1)\cdot 4c\log(inw/\gamma)+(cj\log(1/\gamma)+2c\log(knw/\gamma))\\
&\le ck\log(1/\gamma) +h\cdot 4c\log(knw/\gamma)\\
&\le \si[,h,k]
\end{align*}
which is also true. Moreover, observe that when $k\le 1$ it is always the case that $i,j\le\ceil{k/2}$. Therefore the third condition is also true. Finally we show that the bound for $\sout[,h,k]$ is also correct. First observe that the restriction $\sout[,h-1,i]\le \sout[,h,k]$ is trivially true. Then the only condition left is that for every $i\le\ceil{k/2}$,
$$\sout[,h-1,i]+\si[,h-1,i]+c\log(1/\delta)+c\log(1/\eps_i)\le \sout[,h,k].$$ Since $\sout[,h-1,i]\le \sout[,h-1,\ceil{k/2}]$ and $\si[,h-1,i]\le \si[,h-1,\ceil{k/2}]$ for every $i$, it suffices to show that 
$$\sout[,h-1,\ceil{k/2}]+\si[,h-1,\ceil{k/2}]+c\log(1/\delta)+c\log(1/\eps_{\ceil{k/2}})\le \sout[,h,k].$$
First we consider $k\le 1$, which is the case that $\ceil{k/2}=k$. Then 
\begin{align*}
&\sout[,h-1,\ceil{k/2}]+\si[,h-1,\ceil{k/2}]+c\log(1/\delta)+c\log(1/\eps_{\ceil{k/2}})\\
&\le \sout(\cG_{a,h-1,k}) + 3ck\log(n/\gamma) + 7c\log(n/\gamma)\\
&\le h \cdot (3ck\log(n/\gamma) + 7c\log(n/\gamma)) \\
&\le  \sout[,h,k].
\end{align*}
Finally we consider the case $k>1$. Observe that 
\begin{align*}
&\sout[,h-1,\ceil{k/2}]+\si[,h-1,\ceil{k/2}]+c\log(1/\delta)+c\log(1/\eps_{\ceil{k/2}})\\
&\le \sout[,h-1,\ceil{k/2}]+\si[,h-1,\ceil{k/2}] + \frac{3k+1}{2}\cdot c\log(n/\gamma)+3c\log(w/\gamma)\\
&\le \sout[,h-1,\ceil{k/2}]+ (2k+1)\cdot c\log(n/\gamma) + (h-1)\cdot 4c\log(w/\gamma) + 7c\log(w/\gamma) \\
&\le  4c\cdot\frac{k+1}{2}\cdot\log(n/\gamma) + \left(\ceil{\log \ceil{\frac{k}{2}}}+1\right)\cdot (h-1) \cdot (10c\log(nw/\gamma)) \\
&+ (2k+1)\cdot c\log(n/\gamma) + (h-1)\cdot 4c\log(w/\gamma) + 7c\log(w/\gamma) \\
&\le 4ck\log(n/\gamma) + \left(\ceil{\log \ceil{\frac{k}{2}}}+1\right)\cdot (h-1) \cdot (10c\log(nw/\gamma)) + h\cdot 10c\log(nw/\gamma)  \\
&\le \sout[,h,k].
\end{align*}
In the last inequality we use the fact that $\ceil{\log k}=\ceil{\log(\ceil{k/2})}+1$ for every $k>1$.

Finally, note that $(11^{\log n}\gamma)=n^{\log_2 11}\cdot n^{-4}\le n^{-0.5}$. By taking $h=\log n$ and  $k=\frac{\log(1/\eps)}{\log (1/n^{0.5})}$, we get a $(n,w,\eps)$-robust PRPD.
\end{proof}

\begin{remark}
To get the seed length we claimed in \Cref{thm:main}, observe that the $\log(1/\eps)$ term is dominating when $\log(1/\eps)\ge \log^3(nw)$. Therefore we can simply replace the $\log\log(1/\eps)$ factor on the $O(\log n\log (nw))$ term with $\log\log(nw)$.
\end{remark}

\appendix
\setcounter{section}{1}
\section{Proof of \Cref{lemma:samp-opt}}\label{sec:proof-of-samp}
\begin{lemma}[\Cref{lemma:samp-opt}, restated. \cite{RVW01,Gol11}]
For every $\delta,\eps>0$ and integer $m$, there exists a $(\eps,\delta)$-sampler $f:\bits{n}\times\bits{d}\to\bits{m}$ s.t. $d=O(\log\log(1/\delta)+\log(1/\eps))$ and $n=m+O(\log(1/\delta))+O(\log(1/\eps))$. Moreover, for every $x,y$, $f(x,y)$ can be computed in space $O(m+\log(1/\delta)+\log(1/\eps))$.
\end{lemma}

We will use the equivalence between seeded randomness extractor and oblivious sampler by Zuckerman~\cite{Zuc97}. To achieve the parameter we need, we need a ``high-entropy seeded extractor" such that the seed length only depends on entropy loss but not the length of source. We will use the standard ``block-source" construction for high-entropy extractor which can be found in \cite{GW97,RVW01}. For simplicity, we will use simple composition instead of zig-zag composition~\cite{RVW01} because we are not aiming for optimal entropy loss. We will use the following standard lemmas for the extractor construction. Some of the following lemmas are implicit in their original source, and we recommend the readers to see \cite{GUV09,Vad12} for a proof.
\begin{definition}[\cite{CG88}]
$(X_1,X_2)$ is a \emph{$(k_1,k_2)$-block source} if $X_1$ is a $k_1$-source, and for every $x_1\in\Supp(X)$, $X_2$ conditioned on $X_1=x_1$  is a $k_2$-source.
\end{definition}
\begin{lemma}[\cite{GW97}]\label{lemma:block}
Let $X\in\bits{n}$ be a $(n-\Delta)$ source. Then for every integer $0\le t\le n $, $X$ is $\eps$-close to a $(t-\Delta,n-t-\Delta-\log(1/\eps))$-block source $(X_1,X_2)$ where $X_1\in\bits{t}$ and $X_2\in\bits{n-t}$.
\end{lemma}
\begin{lemma}[\cite{NZ96}]\label{lemma:compo}
Let $E_1:\bits{n_1}\times\bits{d}\to\bits{d_2}$ be a $(k_1,\eps_1)$ extractor and $E_2:\bits{n_2}\times\bits{d_2}\to\bits{m}$ be a $(k_2,\eps_2)$ extractor. Define $E((x_1,x_2),s)=E_1(x_2,E_2(x_1,s))$. Then for every $(k_1,k_2)$-block source $(X_1,X_2)\in\bits{n_1}\times\bits{n_2}$, $E((X_1,X_2),U_d)$ is ${(\eps_1+\eps_2)}$-close to uniform.
\end{lemma}
\begin{lemma}[\cite{GW97}]\label{lemma:GW}
For every $\eps,\Delta>0$ and integer $n$ there exists a $(n-\Delta,\eps)$ extractor $E:\bits{n}\times\bits{d}\to\bits{n}$ with $d=O(\Delta+\log(1/\eps))$, and for every $x,y$, $E(x,y)$ can be computed in space $O(n+\log(1/\eps))$.
\end{lemma}
\begin{lemma}[\cite{GUV09,KNW08}]\label{lemma:GUV}
For every $\eps>0$, integer $m>0$ and $n\ge 2m$, there exists a $(2m,\eps)$ extractor $E:\bits{n}\times\bits{d}\to\bits{m}$ with $d=O(\log m + \log(1/\eps))$, and for every $x,y$, $E(x,y)$  can be computed in space $O(m+\log(1/\eps))$.
\end{lemma}
\begin{lemma}[\cite{Zuc97}]\label{lemma:Zuc}
Every $(n-\log(1/\delta)-1,\eps)$-extractor is a $(\eps,\delta)$-sampler.
\end{lemma}
\noindent
Now we show how to construct the sampler we need, and that it is indeed space efficient.
\begin{proof}
Let $\Delta=\log(1/\delta)+1$. Let $E_1:\bits{m}\times\bits{d_1}\to\bits{m}$ be an $(m-\Delta,\eps/3)$-extractor from Lemma~\ref{lemma:GW}, w.l.o.g. assume that $d_1\ge \Delta+\log(3/\eps)$. Then let $E_2:\bits{3d_1}\times\bits{d}\to\bits{d_1}$ be an $(2d_1,\eps/3)$-extractor from Lemma~\ref{lemma:GUV}. Then we claim that $E:\bits{m+3d_1}\times\bits{d}\to\bits{m}$, defined as $E((x_1,x_2),s)=E_1(x_1,E_2(x_2,s))$, is a $(m+3d_1-\Delta,\eps)$ extractor, and hence a $(\eps,\delta)$ sampler by Lemma~\ref{lemma:Zuc}. 

To prove the claim, consider any $(m+3d_1-\Delta)$-source $X$. By Lemma~\ref{lemma:block}, $X$ is $(\eps/3)$-close to a $(m-\Delta,3d_1-\Delta-\log(3/\eps))$-block source $(X_1,X_2)\in\bits{3d_1}\times\bits{m}$. By Lemma~\ref{lemma:compo}, $E_1(X_1,E_2(X_2,U_d))$ is $2\eps/3$-close to uniform. Since $E(X,U_d)$ is $\eps/3$-close to $E_1(X_1,E_2(X_2,U_d))$, by triangle inequality it is $\eps$-close to uniform. Moreover, $d=O(\log (d_1/\eps))=O(\log\log(1/\delta)+\log(1/\eps))$, $n=m+3d_1=m+O(\log(1/\delta)+\log(1/\eps))$, and the required space to compute $E$ is $O(m+d_1+\log(1/\eps))=O(m+\log(1/\eps)+\log(1/\delta))$.
\end{proof}

\begin{thebibliography}{99}
\bibitem{Arm98}
Roy Armoni.
\newblock On the derandomization of space-bounded computations.
\newblock In Michael Luby, Jos{\'{e}} D.~P. Rolim, and Maria~J. Serna, editors,
  {\em Randomization and Approximation Techniques in Computer Science, Second
  International Workshop, RANDOM'98, Barcelona, Spain, October 8-10, 1998,
  Proceedings}, volume 1518 of {\em Lecture Notes in Computer Science}, pages
  47--59. Springer, 1998.

\bibitem{BCG18}
Mark Braverman, Gil Cohen, and Sumegha Garg.
\newblock Hitting sets with near-optimal error for read-once branching
  programs.
\newblock In Ilias Diakonikolas, David Kempe, and Monika Henzinger, editors,
  {\em Proceedings of the 50th Annual {ACM} {SIGACT} Symposium on Theory of
  Computing, {STOC} 2018, Los Angeles, CA, USA, June 25-29, 2018}, pages
  353--362. {ACM}, 2018.

\bibitem{BDVY13}
Andrej Bogdanov, Zeev Dvir, Elad Verbin, and Amir Yehudayoff.
\newblock Pseudorandomness for width-2 branching programs.
\newblock {\em Theory of Computing}, 9:283--293, 2013.

\bibitem{BRRY14}
Mark Braverman, Anup Rao, Ran Raz, and Amir Yehudayoff.
\newblock Pseudorandom generators for regular branching programs.
\newblock {\em {SIAM} J. Comput.}, 43(3):973--986, 2014.

\bibitem{BV10}
Joshua Brody and Elad Verbin.
\newblock The coin problem and pseudorandomness for branching programs.
\newblock In {\em 51th Annual {IEEE} Symposium on Foundations of Computer
  Science, {FOCS} 2010, October 23-26, 2010, Las Vegas, Nevada, {USA}}, pages
  30--39. {IEEE} Computer Society, 2010.

\bibitem{CG88}
Benny Chor and Oded Goldreich.
\newblock Unbiased bits from sources of weak randomness and probabilistic
  communication complexity.
\newblock {\em {SIAM} J. Comput.}, 17(2):230--261, 1988.

\bibitem{De11}
Anindya De.
\newblock Pseudorandomness for permutation and regular branching programs.
\newblock In {\em Proceedings of the 26th Annual {IEEE} Conference on
  Computational Complexity, {CCC} 2011, San Jose, California, USA, June 8-10,
  2011}, pages 221--231. {IEEE} Computer Society, 2011.

\bibitem{Gol11}
Oded Goldreich.
\newblock A sample of samplers: {A} computational perspective on sampling.
\newblock In Oded Goldreich, editor, {\em Studies in Complexity and
  Cryptography. Miscellanea on the Interplay between Randomness and Computation
  - In Collaboration with Lidor Avigad, Mihir Bellare, Zvika Brakerski, Shafi
  Goldwasser, Shai Halevi, Tali Kaufman, Leonid Levin, Noam Nisan, Dana Ron,
  Madhu Sudan, Luca Trevisan, Salil Vadhan, Avi Wigderson, David Zuckerman},
  volume 6650 of {\em Lecture Notes in Computer Science}, pages 302--332.
  Springer, 2011.

\bibitem{GUV09}
Venkatesan Guruswami, Christopher Umans, and Salil~P. Vadhan.
\newblock Unbalanced expanders and randomness extractors from parvaresh-vardy
  codes.
\newblock {\em J. {ACM}}, 56(4):20:1--20:34, 2009.

\bibitem{GW97}
Oded Goldreich and Avi Wigderson.
\newblock Tiny families of functions with random properties: {A} quality-size
  trade-off for hashing.
\newblock {\em Random Struct. Algorithms}, 11(4):315--343, 1997.

\bibitem{INW94}
Russell Impagliazzo, Noam Nisan, and Avi Wigderson.
\newblock Pseudorandomness for network algorithms.
\newblock In Frank~Thomson Leighton and Michael~T. Goodrich, editors, {\em
  Proceedings of the Twenty-Sixth Annual {ACM} Symposium on Theory of
  Computing, 23-25 May 1994, Montr{\'{e}}al, Qu{\'{e}}bec, Canada}, pages
  356--364. {ACM}, 1994.

\bibitem{KNP11}
Michal Kouck{\'{y}}, Prajakta Nimbhorkar, and Pavel Pudl{\'{a}}k.
\newblock Pseudorandom generators for group products: extended abstract.
\newblock In Lance Fortnow and Salil~P. Vadhan, editors, {\em Proceedings of
  the 43rd {ACM} Symposium on Theory of Computing, {STOC} 2011, San Jose, CA,
  USA, 6-8 June 2011}, pages 263--272. {ACM}, 2011.

\bibitem{KNW08}
Daniel~M. Kane, Jelani Nelson, and David~P. Woodruff.
\newblock Revisiting norm estimation in data streams.
\newblock {\em CoRR}, abs/0811.3648, 2008.

\bibitem{MRT19}
Raghu Meka, Omer Reingold, and Avishay Tal.
\newblock Pseudorandom generators for width-3 branching programs.
\newblock In Moses Charikar and Edith Cohen, editors, {\em Proceedings of the
  51st Annual {ACM} {SIGACT} Symposium on Theory of Computing, {STOC} 2019,
  Phoenix, AZ, USA, June 23-26, 2019}, pages 626--637. {ACM}, 2019.

\bibitem{Nis92}
Noam Nisan.
\newblock Pseudorandom generators for space-bounded computation.
\newblock {\em Combinatorica}, 12(4):449--461, 1992.

\bibitem{NZ96}
Noam Nisan and David Zuckerman.
\newblock Randomness is linear in space.
\newblock {\em J. Comput. Syst. Sci.}, 52(1):43--52, 1996.

\bibitem{RR99}
Ran Raz and Omer Reingold.
\newblock On recycling the randomness of states in space bounded computation.
\newblock In Jeffrey~Scott Vitter, Lawrence~L. Larmore, and Frank~Thomson
  Leighton, editors, {\em Proceedings of the Thirty-First Annual {ACM}
  Symposium on Theory of Computing, May 1-4, 1999, Atlanta, Georgia, {USA}},
  pages 159--168. {ACM}, 1999.

\bibitem{RVW01}
Omer Reingold, Salil~P. Vadhan, and Avi Wigderson.
\newblock Entropy waves, the zig-zag graph product, and new constant-degree
  expanders and extractors.
\newblock {\em Electronic Colloquium on Computational Complexity {(ECCC)}},
  8(18), 2001.

\bibitem{Ste12}
Thomas Steinke.
\newblock Pseudorandomness for permutation branching programs without the group
  theory.
\newblock {\em Electronic Colloquium on Computational Complexity {(ECCC)}},
  19:83, 2012.

\bibitem{Vad12}
Salil~P. Vadhan.
\newblock Pseudorandomness.
\newblock {\em Foundations and Trends in Theoretical Computer Science},
  7(1-3):1--336, 2012.

\bibitem{Zuc97}
David Zuckerman.
\newblock Randomness-optimal oblivious sampling.
\newblock {\em Random Struct. Algorithms}, 11(4):345--367, 1997.

\end{thebibliography}
\end{document}
